\subsection{MULTIPLICATION OF IDEALS}

Let A be a ring and a and b two-sided ideals of A. The set of elements of the form \(x_{1}y_{1} + \cdots + x_{n}y_{n}\) with \(n \geqslant 0\) , \(x_{i} \in a\) and \(y_{i} \in b\) for \(1 \leqslant i \leqslant n\) , is obviously a two-sided ideal of A, which is denoted by ab and called the product of the two-sided ideals a and b. Under this multiplication, the set of two-sided ideals of A is a monoid with unit element the two-sided ideal A. If a, b, c are two-sided ideals of A, then \(a(b + c) = ab + ac\) , \((b + c)a = ba + ca\) . If A is commutative, multiplication of ideals is commutative.

\(\mathsf{ab}\subset \mathsf{aA}\subset \mathsf{a}\) and \(\mathsf{ab}\subset \mathsf{Ab}\subset \mathsf{b}\) , hence

\[
a b \subset a \cap b.\tag{21}
\]

PROPOSITION 6. Let \(a, b_{1}, \ldots, b_{n}\) be two-sided ideals of A. If \(A = a + b_{i}\) for all \(i\) , then \(A = a + b_{1}b_{2} \ldots b_{n} = a + (b_{1} \cap b_{2} \cap \cdots \cap b_{n})\) .

By (21) it suffices to prove that \(\mathbf{A} = \mathfrak{a} + \mathfrak{b}_1\mathfrak{b}_2\ldots \mathfrak{b}_n\) . By induction, it suffices to consider the case where \(n = 2\) . By hypothesis, there exist \(a, a' \in \mathfrak{a}, b_1 \in \mathfrak{b}_1, b_2 \in \mathfrak{b}_2\) such that \(1 = a + b_1 = a' + b_2\) . Then

\[
1 = a ^ {\prime} + (a + b _ {1}) b _ {2} = \left(a ^ {\prime} + a b _ {2}\right) + b _ {1} b _ {2} \in \mathfrak {a} + \mathfrak {b} _ {1} \mathfrak {b} _ {2},
\]

whence \(\mathbf{A} = \mathfrak{a} + \mathfrak{b}_1\mathfrak{b}_2\)

PROPOSITION 7. Let \(\mathfrak{b}_1, \ldots, \mathfrak{b}_n\) be two-sided ideals of A such that \(\mathfrak{b}_i + \mathfrak{b}_j = \mathbf{A}\) for \(i \neq j\) . Then \(\mathfrak{b}_1 \cap \mathfrak{b}_2 \cap \dots \cap \mathfrak{b}_n = \sum_{\sigma \in \mathfrak{E}_n} \mathfrak{b}_{\sigma(1)} \mathfrak{b}_{\sigma(2)} \ldots \mathfrak{b}_{\sigma(n)}\) . In particular, if A is commutative, \(\mathfrak{b}_1 \cap \mathfrak{b}_2 \cap \dots \cap \mathfrak{b}_n = \mathfrak{b}_1 \mathfrak{b}_2 \ldots \mathfrak{b}_n\) (cf.~Exercise 2).

Suppose first \(n = 2\) . There exist \(a_1 \in \mathfrak{b}_1\) , \(a_2 \in \mathfrak{b}_2\) such that \(a_1 + a_2 = 1\) . If \(x \in \mathfrak{b}_1 \cap \mathfrak{b}_2\) , then \(x = x(a_1 + a_2) = xa_1 + xa_2 \in \mathfrak{b}_2\mathfrak{b}_1 + \mathfrak{b}_1\mathfrak{b}_2\) . Hence \(\mathfrak{b}_1 \cap \mathfrak{b}_2 = \mathfrak{b}_1\mathfrak{b}_2 + \mathfrak{b}_2\mathfrak{b}_1\) .

Suppose now the equation of the proposition is established for all integers \(<n\) . By Proposition 6, \(\mathfrak{b}_{n} + (\mathfrak{b}_{1}\mathfrak{b}_{2}\ldots \mathfrak{b}_{n - 1}) = A\) and hence

\[
\begin{array}{l} \mathfrak {b} _ {1} \cap \mathfrak {b} _ {2} \cap \dots \cap \mathfrak {b} _ {n} = (\mathfrak {b} _ {1} \cap \mathfrak {b} _ {2} \cap \dots \cap \mathfrak {b} _ {n - 1}) \mathfrak {b} _ {n} + \mathfrak {b} _ {n} (\mathfrak {b} _ {1} \cap \mathfrak {b} _ {2} \cap \dots \cap \mathfrak {b} _ {n - 1}) \\ = \Big (\sum_ {\sigma \in \mathfrak {E} _ {n - 1}} \mathfrak {b} _ {\sigma (1)} \mathfrak {b} _ {\sigma (2)} \dots \mathfrak {b} _ {\sigma (n - 1)} \Big) \mathfrak {b} _ {n} \\ \qquad + \mathfrak {b} _ {n} \Big (\sum_ {\sigma \in \mathfrak {E} _ {n - 1}} \mathfrak {b} _ {\sigma (1)} \mathfrak {b} _ {\sigma (2)} \dots \mathfrak {b} _ {\sigma (n - 1)} \Big) \\ \subset \sum_ {\tau \in \mathfrak {E} _ {n}} \mathfrak {b} _ {\tau (1)} \mathfrak {b} _ {\tau (2)} \dots \mathfrak {b} _ {\tau (n)} \subset \mathfrak {b} _ {1} \cap \mathfrak {b} _ {2} \cap \dots \cap \mathfrak {b} _ {n}. \end{array}
\]

